% Options for packages loaded elsewhere
\PassOptionsToPackage{unicode}{hyperref}
\PassOptionsToPackage{hyphens}{url}
\PassOptionsToPackage{dvipsnames,svgnames,x11names}{xcolor}
\PassOptionsToPackage{space}{xeCJK}
\documentclass[
  chinese,
  11pt,
  a4paper,
]{article}
\usepackage{xcolor}
\usepackage[margin=2.2cm]{geometry}
\usepackage{amsmath,amssymb}
\setcounter{secnumdepth}{-\maxdimen} % remove section numbering
\usepackage{iftex}
\ifPDFTeX
  \usepackage[T1]{fontenc}
  \usepackage[utf8]{inputenc}
  \usepackage{textcomp} % provide euro and other symbols
\else % if luatex or xetex
  \usepackage{unicode-math} % this also loads fontspec
  \defaultfontfeatures{Scale=MatchLowercase}
  \defaultfontfeatures[\rmfamily]{Ligatures=TeX,Scale=1}
\fi
\usepackage{lmodern}
\ifPDFTeX\else
  % xetex/luatex font selection
  \setmainfont[]{Noto Serif CJK SC}
  \setsansfont[]{Noto Sans CJK SC}
  \setmonofont[]{Noto Sans Mono CJK SC}
  \ifXeTeX
    \usepackage{xeCJK}
    \setCJKmainfont[]{Noto Serif CJK SC}
  \fi
  \ifLuaTeX
    \usepackage[]{luatexja-fontspec}
    \setmainjfont[]{Noto Serif CJK SC}
  \fi
\fi
% Use upquote if available, for straight quotes in verbatim environments
\IfFileExists{upquote.sty}{\usepackage{upquote}}{}
\IfFileExists{microtype.sty}{% use microtype if available
  \usepackage[]{microtype}
  \UseMicrotypeSet[protrusion]{basicmath} % disable protrusion for tt fonts
}{}
\makeatletter
\@ifundefined{KOMAClassName}{% if non-KOMA class
  \IfFileExists{parskip.sty}{%
    \usepackage{parskip}
  }{% else
    \setlength{\parindent}{0pt}
    \setlength{\parskip}{6pt plus 2pt minus 1pt}}
}{% if KOMA class
  \KOMAoptions{parskip=half}}
\makeatother
\usepackage{color}
\usepackage{fancyvrb}
\newcommand{\VerbBar}{|}
\newcommand{\VERB}{\Verb[commandchars=\\\{\}]}
\DefineVerbatimEnvironment{Highlighting}{Verbatim}{commandchars=\\\{\}}
% Add ',fontsize=\small' for more characters per line
\usepackage{framed}
\definecolor{shadecolor}{RGB}{248,248,248}
\newenvironment{Shaded}{\begin{snugshade}}{\end{snugshade}}
\newcommand{\AlertTok}[1]{\textcolor[rgb]{0.94,0.16,0.16}{#1}}
\newcommand{\AnnotationTok}[1]{\textcolor[rgb]{0.56,0.35,0.01}{\textbf{\textit{#1}}}}
\newcommand{\AttributeTok}[1]{\textcolor[rgb]{0.13,0.29,0.53}{#1}}
\newcommand{\BaseNTok}[1]{\textcolor[rgb]{0.00,0.00,0.81}{#1}}
\newcommand{\BuiltInTok}[1]{#1}
\newcommand{\CharTok}[1]{\textcolor[rgb]{0.31,0.60,0.02}{#1}}
\newcommand{\CommentTok}[1]{\textcolor[rgb]{0.56,0.35,0.01}{\textit{#1}}}
\newcommand{\CommentVarTok}[1]{\textcolor[rgb]{0.56,0.35,0.01}{\textbf{\textit{#1}}}}
\newcommand{\ConstantTok}[1]{\textcolor[rgb]{0.56,0.35,0.01}{#1}}
\newcommand{\ControlFlowTok}[1]{\textcolor[rgb]{0.13,0.29,0.53}{\textbf{#1}}}
\newcommand{\DataTypeTok}[1]{\textcolor[rgb]{0.13,0.29,0.53}{#1}}
\newcommand{\DecValTok}[1]{\textcolor[rgb]{0.00,0.00,0.81}{#1}}
\newcommand{\DocumentationTok}[1]{\textcolor[rgb]{0.56,0.35,0.01}{\textbf{\textit{#1}}}}
\newcommand{\ErrorTok}[1]{\textcolor[rgb]{0.64,0.00,0.00}{\textbf{#1}}}
\newcommand{\ExtensionTok}[1]{#1}
\newcommand{\FloatTok}[1]{\textcolor[rgb]{0.00,0.00,0.81}{#1}}
\newcommand{\FunctionTok}[1]{\textcolor[rgb]{0.13,0.29,0.53}{\textbf{#1}}}
\newcommand{\ImportTok}[1]{#1}
\newcommand{\InformationTok}[1]{\textcolor[rgb]{0.56,0.35,0.01}{\textbf{\textit{#1}}}}
\newcommand{\KeywordTok}[1]{\textcolor[rgb]{0.13,0.29,0.53}{\textbf{#1}}}
\newcommand{\NormalTok}[1]{#1}
\newcommand{\OperatorTok}[1]{\textcolor[rgb]{0.81,0.36,0.00}{\textbf{#1}}}
\newcommand{\OtherTok}[1]{\textcolor[rgb]{0.56,0.35,0.01}{#1}}
\newcommand{\PreprocessorTok}[1]{\textcolor[rgb]{0.56,0.35,0.01}{\textit{#1}}}
\newcommand{\RegionMarkerTok}[1]{#1}
\newcommand{\SpecialCharTok}[1]{\textcolor[rgb]{0.81,0.36,0.00}{\textbf{#1}}}
\newcommand{\SpecialStringTok}[1]{\textcolor[rgb]{0.31,0.60,0.02}{#1}}
\newcommand{\StringTok}[1]{\textcolor[rgb]{0.31,0.60,0.02}{#1}}
\newcommand{\VariableTok}[1]{\textcolor[rgb]{0.00,0.00,0.00}{#1}}
\newcommand{\VerbatimStringTok}[1]{\textcolor[rgb]{0.31,0.60,0.02}{#1}}
\newcommand{\WarningTok}[1]{\textcolor[rgb]{0.56,0.35,0.01}{\textbf{\textit{#1}}}}
\ifLuaTeX
\usepackage[bidi=basic,shorthands=off,provide=*]{babel}
\else
\usepackage[bidi=default,shorthands=off,provide=*]{babel}
\fi
\ifPDFTeX
\else
\babelfont{rm}[]{Noto Serif CJK SC}
\fi
\ifLuaTeX
  \usepackage{selnolig} % disable illegal ligatures
\fi
\setlength{\emergencystretch}{3em} % prevent overfull lines
\providecommand{\tightlist}{%
  \setlength{\itemsep}{0pt}\setlength{\parskip}{0pt}}
\usepackage{bookmark}
\IfFileExists{xurl.sty}{\usepackage{xurl}}{} % add URL line breaks if available
\urlstyle{same}
% fallback for those not using the hyperref driver hyperxmp:
\makeatletter
\@ifundefined{xmpquote}{\newcommand{\xmpquote}[1]{#1}}{}
\makeatother
\hypersetup{
  pdftitle={RGB-D 书本厚度与夹爪开口计算------客户操作手册},
  pdflang={zh-CN},
  colorlinks=true,
  linkcolor={blue},
  filecolor={Maroon},
  citecolor={Blue},
  urlcolor={blue},
  pdfcreator={LaTeX via pandoc}}

\title{RGB-D 书本厚度与夹爪开口计算------客户操作手册}
\author{}
\date{}

\begin{document}
\maketitle

\renewcommand*\contentsname{目录}
{
\setcounter{tocdepth}{3}
\tableofcontents
}
需要了解算法公式、平面分解、侧面剔除和安全拒绝条件时，请查看
\href{BOOK_THICKNESS_ALGORITHM.md}{\texttt{BOOK\_THICKNESS\_ALGORITHM.md}}。

本手册用于从通用 RGB-D 文件或 Record3D
扫描数据中自动识别每一本书，计算书脊厚度，
并生成双指夹爪的建议预张开宽度。

系统输出使用以下公式：

\begin{Shaded}
\begin{Highlighting}[]
\NormalTok{建议夹爪开口 = 书本厚度 + 95\% 测量不确定度 + 安全间隙}
\end{Highlighting}
\end{Shaded}

默认安全间隙为 3 mm。夹爪接触书本后，仍应使用力控或电流阈值完成闭合。

\subsection{1. 设置路径}\label{1-ux8bbeux7f6eux8defux5f84}

打开终端，依次执行：

\begin{Shaded}
\begin{Highlighting}[]
\BuiltInTok{export} \VariableTok{PROJECT}\OperatorTok{=}\NormalTok{/mnt/ssd/user/V/feature{-}splatting{-}inria}
\BuiltInTok{export} \VariableTok{CAPTURE}\OperatorTok{=}\VariableTok{$PROJECT}\NormalTok{/2026{-}09{-}01{-}{-}23{-}22{-}25}
\BuiltInTok{export} \VariableTok{PYTHON}\OperatorTok{=}\NormalTok{/mnt/ssd/robotarm/env/sam3\_full/bin/python}
\end{Highlighting}
\end{Shaded}

命令说明：

\begin{itemize}
\tightlist
\item
  \texttt{PROJECT}：厚度计算程序所在目录。
\item
  \texttt{CAPTURE}：本次 Record3D 扫描数据目录。
\item
  \texttt{PYTHON}：已经安装 SAM3 和 PyTorch 的 Python 解释器。
\end{itemize}

如需分析其他扫描，只需修改 \texttt{CAPTURE}：

\begin{Shaded}
\begin{Highlighting}[]
\BuiltInTok{export} \VariableTok{CAPTURE}\OperatorTok{=}\NormalTok{/path/to/another{-}record3d{-}capture}
\end{Highlighting}
\end{Shaded}

\subsection{2.
检查扫描数据}\label{2-ux68c0ux67e5ux626bux63cfux6570ux636e}

执行：

\begin{Shaded}
\begin{Highlighting}[]
\BuiltInTok{test} \AttributeTok{{-}f} \StringTok{"}\VariableTok{$CAPTURE}\StringTok{/metadata"} \KeywordTok{\&\&} \BuiltInTok{echo} \StringTok{"metadata: OK"}
\BuiltInTok{test} \AttributeTok{{-}f} \StringTok{"}\VariableTok{$CAPTURE}\StringTok{/rgbd/0.jpg"} \KeywordTok{\&\&} \BuiltInTok{echo} \StringTok{"RGB: OK"}
\BuiltInTok{test} \AttributeTok{{-}f} \StringTok{"}\VariableTok{$CAPTURE}\StringTok{/rgbd/0.depth"} \KeywordTok{\&\&} \BuiltInTok{echo} \StringTok{"depth: OK"}
\BuiltInTok{test} \AttributeTok{{-}f} \StringTok{"}\VariableTok{$CAPTURE}\StringTok{/rgbd/0.conf"} \KeywordTok{\&\&} \BuiltInTok{echo} \StringTok{"confidence: OK"}
\end{Highlighting}
\end{Shaded}

本节检查的是 Record3D 输入，多视角模式四项均应显示 \texttt{OK}。如果使用
\texttt{CAPTURE\ -\/-single-frame-rgbd} 兼容模式，只要求
\texttt{metadata}、\texttt{.jpg} 和 \texttt{.depth}。 如果直接使用
\texttt{-\/-rgb}、\texttt{-\/-depth}、\texttt{-\/-intrinsics}
三个通用路径，则不需要 Record3D 目录和 \texttt{metadata}，请直接查看 6.1
节。Record3D 数据目录的基本结构为：

\begin{Shaded}
\begin{Highlighting}[]
\NormalTok{CAPTURE/}
\NormalTok{├── metadata}
\NormalTok{└── rgbd/}
\NormalTok{    ├── 0.jpg}
\NormalTok{    ├── 0.depth}
\NormalTok{    ├── 0.conf}
\NormalTok{    ├── 1.jpg}
\NormalTok{    └── ...}
\end{Highlighting}
\end{Shaded}

其中：

\begin{itemize}
\tightlist
\item
  \texttt{.jpg} 是高分辨率 RGB 图像。
\item
  \texttt{.depth} 是 Record3D LiDAR 深度。
\item
  \texttt{.conf} 是深度置信度；单帧 RGB-D 模式可省略。
\item
  \texttt{metadata}
  必须包含相机内参和图像尺寸；只有多视角模式要求其中包含位姿。
\end{itemize}

\subsection{3. 检查 Python 和
GPU}\label{3-ux68c0ux67e5-python-ux548c-gpu}

检查 Python：

\begin{Shaded}
\begin{Highlighting}[]
\StringTok{"}\VariableTok{$PYTHON}\StringTok{"} \AttributeTok{{-}{-}version}
\end{Highlighting}
\end{Shaded}

检查 PyTorch 与 GPU：

\begin{Shaded}
\begin{Highlighting}[]
\StringTok{"}\VariableTok{$PYTHON}\StringTok{"} \AttributeTok{{-}c} \StringTok{"import torch; print(\textquotesingle{}PyTorch:\textquotesingle{}, torch.\_\_version\_\_); print(\textquotesingle{}CUDA:\textquotesingle{}, torch.cuda.is\_available()); print(\textquotesingle{}GPU:\textquotesingle{}, torch.cuda.get\_device\_name(0) if torch.cuda.is\_available() else \textquotesingle{}不可用\textquotesingle{})"}
\end{Highlighting}
\end{Shaded}

正常情况下应看到：

\begin{Shaded}
\begin{Highlighting}[]
\NormalTok{CUDA: True}
\NormalTok{GPU: NVIDIA ...}
\end{Highlighting}
\end{Shaded}

如果显示 \texttt{CUDA:\ False}，请先检查 NVIDIA 驱动：

\begin{Shaded}
\begin{Highlighting}[]
\ExtensionTok{nvidia{-}smi}
\end{Highlighting}
\end{Shaded}

\subsection{4. 一次性安装 Record3D
解码依赖}\label{4-ux4e00ux6b21ux6027ux5b89ux88c5-record3d-ux89e3ux7801ux4f9dux8d56}

首次部署时执行：

\begin{Shaded}
\begin{Highlighting}[]
\StringTok{"}\VariableTok{$PYTHON}\StringTok{"} \AttributeTok{{-}m}\NormalTok{ pip install pyliblzfse}
\end{Highlighting}
\end{Shaded}

验证安装：

\begin{Shaded}
\begin{Highlighting}[]
\StringTok{"}\VariableTok{$PYTHON}\StringTok{"} \AttributeTok{{-}c} \StringTok{"import liblzfse; print(\textquotesingle{}Record3D decoder: OK\textquotesingle{})"}
\end{Highlighting}
\end{Shaded}

看到 \texttt{Record3D\ decoder:\ OK}
即表示安装成功。此步骤只需执行一次。

\subsection{5. 检查 SAM3 模型}\label{5-ux68c0ux67e5-sam3-ux6a21ux578b}

当前部署使用以下模型：

\begin{Shaded}
\begin{Highlighting}[]
\BuiltInTok{export} \VariableTok{SAM3\_CHECKPOINT}\OperatorTok{=}\NormalTok{/home/vlmdemo/.cache/huggingface/hub/models{-}{-}facebook{-}{-}sam3/snapshots/3c879f39826c281e95690f02c7821c4de09afae7/sam3.pt}
\BuiltInTok{test} \AttributeTok{{-}f} \StringTok{"}\VariableTok{$SAM3\_CHECKPOINT}\StringTok{"} \KeywordTok{\&\&} \BuiltInTok{echo} \StringTok{"SAM3 checkpoint: OK"}
\end{Highlighting}
\end{Shaded}

如果没有显示 \texttt{OK}，需要将 SAM3 权重放到该位置，或在运行时通过
\texttt{-\/-checkpoint} 指定实际路径。

\subsection{6.
使用默认参数运行}\label{6-ux4f7fux7528ux9ed8ux8ba4ux53c2ux6570ux8fd0ux884c}

进入项目目录：

\begin{Shaded}
\begin{Highlighting}[]
\BuiltInTok{cd} \StringTok{"}\VariableTok{$PROJECT}\StringTok{"}
\end{Highlighting}
\end{Shaded}

执行完整分析：

\begin{Shaded}
\begin{Highlighting}[]
\VariableTok{CUDA\_VISIBLE\_DEVICES}\OperatorTok{=}\NormalTok{0 }\DataTypeTok{\textbackslash{}}
\StringTok{"}\VariableTok{$PYTHON}\StringTok{"}\NormalTok{ book\_thickness\_pipeline.py }\StringTok{"}\VariableTok{$CAPTURE}\StringTok{"}
\end{Highlighting}
\end{Shaded}

逐行说明：

\begin{itemize}
\tightlist
\item
  \texttt{CUDA\_VISIBLE\_DEVICES=0}：使用第 1 块 GPU。
\item
  \texttt{"\$PYTHON"}：使用 SAM3 专用 Python 环境。
\item
  \texttt{book\_thickness\_pipeline.py}：运行厚度计算程序。
\item
  \texttt{"\$CAPTURE"}：指定要分析的 RGB-D 扫描。
\end{itemize}

默认行为：

\begin{itemize}
\tightlist
\item
  使用帧 \texttt{0,12,24,36,48,60,71}。
\item
  使用第 24 帧作为边界和编号参考图。
\item
  按每层书架从左到右编号。
\item
  过滤低置信度和重复分割。
\item
  融合多个视角的厚度测量。
\item
  在 95\% 不确定度之外再增加 3 mm 夹爪间隙。
\end{itemize}

正常结束时会显示类似：

\begin{Shaded}
\begin{Highlighting}[]
\NormalTok{Measured 96 books}
\NormalTok{Per shelf: S1=29, S2=33, S3=34}
\NormalTok{Quality: \{\textquotesingle{}high\textquotesingle{}: 91, \textquotesingle{}medium\textquotesingle{}: 5, \textquotesingle{}low\textquotesingle{}: 0\}}
\NormalTok{Results: .../book\_thickness\_analysis}
\end{Highlighting}
\end{Shaded}

\subsubsection{6.1 只用单帧 RGB-D
运行}\label{61-ux53eaux7528ux5355ux5e27-rgb-d-ux8fd0ux884c}

\paragraph{6.1.1
什么时候使用}\label{611-ux4ec0ux4e48ux65f6ux5019ux4f7fux7528}

如果只有一张 RGB
图、一张已配准深度图和一个相机内参文件，可以直接输入这三个
文件的路径。数据不必来自 Record3D，也不要求固定的目录结构。

同时提供 \texttt{-\/-rgb}、\texttt{-\/-depth} 和 \texttt{-\/-intrinsics}
后，程序会自动进入通用单帧模式， 不需要提供
\texttt{CAPTURE}、\texttt{metadata}、\texttt{.conf}、位姿或帧编号。

该模式的行为是：

\begin{itemize}
\tightlist
\item
  直接读取命令行指定的 RGB、depth 和内参文件。
\item
  不依赖 Record3D 的 \texttt{metadata} 和 \texttt{rgbd/} 目录。
\item
  不读取深度 confidence，不使用相机位姿。
\item
  不进行多视角匹配和厚度融合。
\item
  仍使用相机内参将图像像素宽度换算为毫米。
\end{itemize}

原来的 Record3D 单帧命令仍然兼容，详见 6.1.5。

\paragraph{6.1.2
三个输入文件}\label{612-ux4e09ux4e2aux8f93ux5165ux6587ux4ef6}

三个文件可以放在任意目录，例如：

\begin{Shaded}
\begin{Highlighting}[]
\NormalTok{/data/books/rgb.png}
\NormalTok{/data/books/depth.png}
\NormalTok{/data/books/intrinsics.json}
\end{Highlighting}
\end{Shaded}

RGB 支持常用的 JPEG、PNG 和 TIFF。

depth 支持：

\begin{itemize}
\tightlist
\item
  单通道 PNG 或 TIFF。
\item
  \texttt{.npy} 二维 NumPy 数组。
\item
  \texttt{.npz} 文件；数组键建议命名为
  \texttt{depth}。如果只有一个数组也会自动读取。
\end{itemize}

深度单位规则：

\begin{itemize}
\tightlist
\item
  浮点深度默认按米读取。
\item
  16 位及以上整数深度默认按毫米读取。
\item
  如果传感器使用其他单位，必须通过 \texttt{-\/-depth-scale}
  指定每个数值对应多少米。
\end{itemize}

例如，深度图每个单位代表 0.1 mm：

\begin{Shaded}
\begin{Highlighting}[]
\ExtensionTok{{-}{-}depth{-}scale}\NormalTok{ 0.0001}
\end{Highlighting}
\end{Shaded}

推荐将相机内参保存为 JSON：

\begin{Shaded}
\begin{Highlighting}[]
\FunctionTok{\{}
  \DataTypeTok{"fx"}\FunctionTok{:} \FloatTok{1348.5952}\FunctionTok{,}
  \DataTypeTok{"fy"}\FunctionTok{:} \FloatTok{1348.5952}\FunctionTok{,}
  \DataTypeTok{"cx"}\FunctionTok{:} \FloatTok{720.9918}\FunctionTok{,}
  \DataTypeTok{"cy"}\FunctionTok{:} \FloatTok{962.8973}
\FunctionTok{\}}
\end{Highlighting}
\end{Shaded}

也可以使用标准的 3×3 相机矩阵：

\begin{Shaded}
\begin{Highlighting}[]
\FunctionTok{\{}
  \DataTypeTok{"K"}\FunctionTok{:} \OtherTok{[}
    \OtherTok{[}\FloatTok{1348.5952}\OtherTok{,} \FloatTok{0.0}\OtherTok{,} \FloatTok{720.9918}\OtherTok{],}
    \OtherTok{[}\FloatTok{0.0}\OtherTok{,} \FloatTok{1348.5952}\OtherTok{,} \FloatTok{962.8973}\OtherTok{],}
    \OtherTok{[}\FloatTok{0.0}\OtherTok{,} \FloatTok{0.0}\OtherTok{,} \FloatTok{1.0}\OtherTok{]}
  \OtherTok{]}
\FunctionTok{\}}
\end{Highlighting}
\end{Shaded}

内参文件还支持 TXT、NPY、NPZ 和 OpenCV YAML/XML。TXT 可以写
\texttt{fx\ fy\ cx\ cy} 四个数，也可以写按行排列的 3×3 矩阵。

内参必须对应 RGB 图像坐标系。不能直接使用未经配准的原始深度相机内参。

RGB 和 depth
必须来自同一次采集并完成空间配准。两者允许分辨率不同，但宽高比
必须接近；程序会拒绝宽高比相差超过 2\% 的输入。

\paragraph{6.1.3
设置输入路径}\label{613-ux8bbeux7f6eux8f93ux5165ux8defux5f84}

执行：

\begin{Shaded}
\begin{Highlighting}[]
\BuiltInTok{export} \VariableTok{RGB}\OperatorTok{=}\NormalTok{/data/books/rgb.png}
\BuiltInTok{export} \VariableTok{DEPTH}\OperatorTok{=}\NormalTok{/data/books/depth.png}
\BuiltInTok{export} \VariableTok{INTRINSICS}\OperatorTok{=}\NormalTok{/data/books/intrinsics.json}
\BuiltInTok{export} \VariableTok{RESULT}\OperatorTok{=}\NormalTok{/data/books/book\_thickness\_result}
\end{Highlighting}
\end{Shaded}

检查文件：

\begin{Shaded}
\begin{Highlighting}[]
\BuiltInTok{test} \AttributeTok{{-}f} \StringTok{"}\VariableTok{$RGB}\StringTok{"} \KeywordTok{\&\&} \BuiltInTok{echo} \StringTok{"RGB: OK"}
\BuiltInTok{test} \AttributeTok{{-}f} \StringTok{"}\VariableTok{$DEPTH}\StringTok{"} \KeywordTok{\&\&} \BuiltInTok{echo} \StringTok{"depth: OK"}
\BuiltInTok{test} \AttributeTok{{-}f} \StringTok{"}\VariableTok{$INTRINSICS}\StringTok{"} \KeywordTok{\&\&} \BuiltInTok{echo} \StringTok{"intrinsics: OK"}
\end{Highlighting}
\end{Shaded}

三项都应显示 \texttt{OK}。

检查 RGB、depth 的尺寸、数据类型和内参：

\begin{Shaded}
\begin{Highlighting}[]
\VariableTok{RGB}\OperatorTok{=}\StringTok{"}\VariableTok{$RGB}\StringTok{"} \VariableTok{DEPTH}\OperatorTok{=}\StringTok{"}\VariableTok{$DEPTH}\StringTok{"} \VariableTok{INTRINSICS}\OperatorTok{=}\StringTok{"}\VariableTok{$INTRINSICS}\StringTok{"} \StringTok{"}\VariableTok{$PYTHON}\StringTok{"} \AttributeTok{{-}} \OperatorTok{\textless{}\textless{}\textquotesingle{}PY\textquotesingle{}}
\StringTok{import json}
\StringTok{import os}
\StringTok{from pathlib import Path}

\StringTok{import cv2}
\StringTok{from PIL import Image}

\StringTok{rgb = Image.open(os.environ["RGB"])}
\StringTok{depth = cv2.imread(os.environ["DEPTH"], cv2.IMREAD\_UNCHANGED)}
\StringTok{intrinsics = json.loads(Path(os.environ["INTRINSICS"]).read\_text())}

\StringTok{print("RGB 尺寸:", rgb.size)}
\StringTok{print("depth 尺寸:", None if depth is None else depth.shape)}
\StringTok{print("depth 类型:", None if depth is None else depth.dtype)}
\StringTok{print("内参:", intrinsics)}
\OperatorTok{PY}
\end{Highlighting}
\end{Shaded}

如果 depth 是 \texttt{.npy} 或 \texttt{.npz}，上面的 OpenCV
检查不适用，可以执行：

\begin{Shaded}
\begin{Highlighting}[]
\VariableTok{DEPTH}\OperatorTok{=}\StringTok{"}\VariableTok{$DEPTH}\StringTok{"} \StringTok{"}\VariableTok{$PYTHON}\StringTok{"} \AttributeTok{{-}} \OperatorTok{\textless{}\textless{}\textquotesingle{}PY\textquotesingle{}}
\StringTok{import os}
\StringTok{import numpy as np}

\StringTok{data = np.load(os.environ["DEPTH"], allow\_pickle=False)}
\StringTok{if isinstance(data, np.lib.npyio.NpzFile):}
\StringTok{    print(\{name: data[name].shape for name in data.files\})}
\StringTok{    data.close()}
\StringTok{else:}
\StringTok{    print("depth 尺寸:", data.shape, "类型:", data.dtype)}
\OperatorTok{PY}
\end{Highlighting}
\end{Shaded}

\paragraph{6.1.4
执行通用单帧分析}\label{614-ux6267ux884cux901aux7528ux5355ux5e27ux5206ux6790}

进入项目目录：

\begin{Shaded}
\begin{Highlighting}[]
\BuiltInTok{cd} \StringTok{"}\VariableTok{$PROJECT}\StringTok{"}
\end{Highlighting}
\end{Shaded}

执行：

\begin{Shaded}
\begin{Highlighting}[]
\VariableTok{CUDA\_VISIBLE\_DEVICES}\OperatorTok{=}\NormalTok{0 }\DataTypeTok{\textbackslash{}}
\StringTok{"}\VariableTok{$PYTHON}\StringTok{"}\NormalTok{ book\_thickness\_pipeline.py }\DataTypeTok{\textbackslash{}}
  \AttributeTok{{-}{-}rgb} \StringTok{"}\VariableTok{$RGB}\StringTok{"} \DataTypeTok{\textbackslash{}}
  \AttributeTok{{-}{-}depth} \StringTok{"}\VariableTok{$DEPTH}\StringTok{"} \DataTypeTok{\textbackslash{}}
  \AttributeTok{{-}{-}intrinsics} \StringTok{"}\VariableTok{$INTRINSICS}\StringTok{"} \DataTypeTok{\textbackslash{}}
  \AttributeTok{{-}{-}output} \StringTok{"}\VariableTok{$RESULT}\StringTok{"}
\end{Highlighting}
\end{Shaded}

逐项说明：

\begin{itemize}
\tightlist
\item
  \texttt{-\/-rgb}：RGB 图像路径。
\item
  \texttt{-\/-depth}：已配准的单通道深度路径。
\item
  \texttt{-\/-intrinsics}：RGB 相机内参路径。
\item
  \texttt{-\/-output}：结果目录；可以省略。
\end{itemize}

如果省略 \texttt{-\/-output}，程序默认在 RGB 所在目录创建
\texttt{RGB文件名\_book\_thickness/}。

如果整数深度不是毫米，增加正确的换算比例。例如每单位为 0.1 mm：

\begin{Shaded}
\begin{Highlighting}[]
\VariableTok{CUDA\_VISIBLE\_DEVICES}\OperatorTok{=}\NormalTok{0 }\DataTypeTok{\textbackslash{}}
\StringTok{"}\VariableTok{$PYTHON}\StringTok{"}\NormalTok{ book\_thickness\_pipeline.py }\DataTypeTok{\textbackslash{}}
  \AttributeTok{{-}{-}rgb} \StringTok{"}\VariableTok{$RGB}\StringTok{"} \DataTypeTok{\textbackslash{}}
  \AttributeTok{{-}{-}depth} \StringTok{"}\VariableTok{$DEPTH}\StringTok{"} \DataTypeTok{\textbackslash{}}
  \AttributeTok{{-}{-}intrinsics} \StringTok{"}\VariableTok{$INTRINSICS}\StringTok{"} \DataTypeTok{\textbackslash{}}
  \AttributeTok{{-}{-}depth{-}scale}\NormalTok{ 0.0001 }\DataTypeTok{\textbackslash{}}
  \AttributeTok{{-}{-}output} \StringTok{"}\VariableTok{$RESULT}\StringTok{"}
\end{Highlighting}
\end{Shaded}

正常结束时会显示类似：

\begin{Shaded}
\begin{Highlighting}[]
\NormalTok{Processing frame 0...}
\NormalTok{  accepted 96 books [29, 33, 34]}
\NormalTok{Measured 96 books}
\NormalTok{Per shelf: S1=29, S2=33, S3=34}
\NormalTok{Quality: \{\textquotesingle{}high\textquotesingle{}: 0, \textquotesingle{}medium\textquotesingle{}: 0, \textquotesingle{}low\textquotesingle{}: 96\}}
\NormalTok{Results: /data/books/book\_thickness\_result}
\end{Highlighting}
\end{Shaded}

\paragraph{6.1.5 Record3D
单帧兼容命令}\label{615-record3d-ux5355ux5e27ux517cux5bb9ux547dux4ee4}

已有 Record3D 目录时，原来的命令仍然可用：

\begin{Shaded}
\begin{Highlighting}[]
\VariableTok{CUDA\_VISIBLE\_DEVICES}\OperatorTok{=}\NormalTok{0 }\DataTypeTok{\textbackslash{}}
\StringTok{"}\VariableTok{$PYTHON}\StringTok{"}\NormalTok{ book\_thickness\_pipeline.py }\StringTok{"}\VariableTok{$CAPTURE}\StringTok{"} \DataTypeTok{\textbackslash{}}
  \AttributeTok{{-}{-}single{-}frame{-}rgbd} \DataTypeTok{\textbackslash{}}
  \AttributeTok{{-}{-}reference{-}frame}\NormalTok{ 24 }\DataTypeTok{\textbackslash{}}
  \AttributeTok{{-}{-}output} \StringTok{"}\VariableTok{$CAPTURE}\StringTok{/book\_thickness\_single\_frame"}
\end{Highlighting}
\end{Shaded}

此兼容模式从 \texttt{CAPTURE/rgbd/24.jpg} 和 \texttt{24.depth}
读取数据，并从 \texttt{metadata} 读取尺寸和内参，但不读取
\texttt{24.conf}，也不使用 \texttt{metadata.poses}。

\paragraph{6.1.6
程序如何计算厚度}\label{616-ux7a0bux5e8fux5982ux4f55ux8ba1ux7b97ux539aux5ea6}

单帧模式按以下顺序处理：

\begin{enumerate}
\def\labelenumi{\arabic{enumi}.}
\tightlist
\item
  在高分辨率 RGB 图像中使用 SAM3 分割每个书脊。
\item
  去除重复分割，并按书架层和左右顺序分配书本 ID。
\item
  根据书脊轮廓拟合长轴和横向轴，得到左右边界；倾斜书本不会直接使用水平框宽。
\item
  将书脊区域映射到低分辨率深度图。
\item
  不使用
  \texttt{.conf}，但会删除无效、过近、过远和明显偏离中位数的深度。
\item
  使用相机内参把深度像素反投影成相机坐标系中的 3D 点。
\item
  每本书使用 RANSAC 分解书脊平面和可能出现的侧面平面。
\item
  汇总画面内所有书本，估计共同的书脊朝向。
\item
  如果掩码边缘存在与书脊方向明显不同的侧面，使用书脊平面的深度支持区裁剪侧面。
\item
  将修正后的左右边界射线与书脊平面求交，计算两个 3D 交点间的毫米距离。
\item
  无法区分书脊与侧面时拒绝该候选，不生成夹爪命令。
\item
  根据通过检查的厚度、不确定度和安全间隙计算夹爪建议开口。
\end{enumerate}

夹爪开口公式为：

\begin{Shaded}
\begin{Highlighting}[]
\NormalTok{建议开口 = 书本厚度 + 95\% 工程不确定度 + 安全间隙}
\end{Highlighting}
\end{Shaded}

默认安全间隙是 3 mm，可以通过 \texttt{-\/-gripper-clearance-mm} 调整。

为避免书本排列松散、斜视时把侧面算进厚度，程序会拒绝以下候选：

\begin{itemize}
\tightlist
\item
  书脊平面深度内点不足。
\item
  书脊平面没有覆盖足够的书脊长度。
\item
  当前书本平面方向与整层书的共同书脊方向差别过大。
\item
  观察角度过于倾斜。
\item
  侧面占比过大。
\item
  已检测到侧面，但无法稳定确定书脊与侧面的交界。
\end{itemize}

被拒绝候选不会进入
\texttt{gripper\_openings.json}，机器人不得为其执行自动抓取。

\paragraph{6.1.7 确认输入来源以及没有使用 confidence 和
pose}\label{617-ux786eux8ba4ux8f93ux5165ux6765ux6e90ux4ee5ux53caux6ca1ux6709ux4f7fux7528-confidence-ux548c-pose}

查看结果前设置：

通用模式已经在 6.1.3 中设置了 \texttt{RESULT}。执行：

\begin{Shaded}
\begin{Highlighting}[]
\VariableTok{RESULT}\OperatorTok{=}\StringTok{"}\VariableTok{$RESULT}\StringTok{"} \StringTok{"}\VariableTok{$PYTHON}\StringTok{"} \AttributeTok{{-}} \OperatorTok{\textless{}\textless{}\textquotesingle{}PY\textquotesingle{}}
\StringTok{import json}
\StringTok{import os}
\StringTok{from pathlib import Path}

\StringTok{result = Path(os.environ["RESULT"])}
\StringTok{data = json.loads((result / "book\_thickness.json").read\_text())}

\StringTok{print("运行模式:", data["mode"])}
\StringTok{print("使用 confidence:", data["uses\_confidence"])}
\StringTok{print("使用相机位姿:", data["uses\_camera\_poses"])}
\StringTok{print("处理帧:", data["frames\_requested"])}
\StringTok{print("输入文件:", data["input\_files"])}
\StringTok{print("书本数量:", data["book\_count"])}
\StringTok{print("拒绝数量:", data["reference\_rejection\_count"])}
\OperatorTok{PY}
\end{Highlighting}
\end{Shaded}

预期结果：

\begin{Shaded}
\begin{Highlighting}[]
\NormalTok{运行模式: single\_frame\_rgbd}
\NormalTok{使用 confidence: False}
\NormalTok{使用相机位姿: False}
\NormalTok{处理帧: [0]}
\NormalTok{输入文件: \{\textquotesingle{}rgb\textquotesingle{}: \textquotesingle{}...\textquotesingle{}, \textquotesingle{}depth\textquotesingle{}: \textquotesingle{}...\textquotesingle{}, \textquotesingle{}intrinsics\textquotesingle{}: \textquotesingle{}...\textquotesingle{}, \textquotesingle{}depth\_scale\_m\_per\_unit\textquotesingle{}: 0.001\}}
\NormalTok{书本数量: \textless{}通过几何检查的书本数\textgreater{}}
\NormalTok{拒绝数量: \textless{}未通过侧面/视角检查的候选数\textgreater{}}
\end{Highlighting}
\end{Shaded}

通用文件没有 Record3D 帧编号，因此程序将本次观测统一记为第 0 帧。
\texttt{input\_files}
会记录实际输入路径以及最终采用的深度比例，便于复查。

\paragraph{6.1.8
如何理解单帧质量}\label{618-ux5982ux4f55ux7406ux89e3ux5355ux5e27ux8d28ux91cf}

单帧无法估计跨视角重复性，因此每本书都会输出：

\begin{Shaded}
\begin{Highlighting}[]
\NormalTok{views\_observed = 1}
\NormalTok{views\_used = 1}
\NormalTok{repeatability\_mad\_mm = 0}
\NormalTok{quality = low}
\end{Highlighting}
\end{Shaded}

注意：

\begin{itemize}
\tightlist
\item
  \texttt{repeatability\_mad\_mm\ =\ 0}
  表示没有第二次观测，不代表测量误差为零。
\item
  \texttt{quality\ =\ low} 是单帧模式的预期行为，不代表程序执行失败。
\item
  \texttt{high\_confidence\_samples}
  不代表真实的传感器高置信度点，因为没有读取 \texttt{.conf}。
\item
  判断是否使用 confidence，应查看顶层的 \texttt{uses\_confidence}。
\item
  单帧 \texttt{uncertainty\_95\_mm}
  不包含真实的跨视角波动，可能低估异常书本的误差。
\item
  \texttt{reference\_side\_face\_detected\ =\ true}：检测到可见侧面。
\item
  \texttt{reference\_boundary\_corrected\ =\ true}：已裁剪侧面，图中书本
  ID 后带 \texttt{*}。
\item
  \texttt{reference\_side\_trim\_mm}：从原轮廓宽度中去除的侧面毫米数。
\item
  \texttt{reference\_view\_angle\_deg}：书脊平面相对正视方向的夹角。
\item
  \texttt{frame\_statistics{[}{]}.rejected\_detections}：被拒绝候选及具体原因。
\end{itemize}

\paragraph{6.1.9
当前数据的测试结果}\label{619-ux5f53ux524dux6570ux636eux7684ux6d4bux8bd5ux7ed3ux679c}

已使用机器人真机采集的以下三个文件完成侧面剔除测试：

\begin{Shaded}
\begin{Highlighting}[]
\NormalTok{dual\_v2\_1/11/rgb/000001.png}
\NormalTok{dual\_v2\_1/11/depth/000001.png}
\NormalTok{dual\_v2\_1/11/cam\_K.txt}
\end{Highlighting}
\end{Shaded}

测试结果：

\begin{itemize}
\tightlist
\item
  SAM3 得到 66 个书本候选。
\item
  57 个候选通过书脊平面几何检查。
\item
  9 个候选因平面方向不一致、书脊内点不足或侧面边界不明确而被拒绝。
\item
  6 本书检测到侧面并完成边界裁剪。
\item
  共从轮廓宽度中裁掉 14.58 mm，单本裁剪 1.86--3.60 mm。
\item
  被拒绝候选没有写入 \texttt{gripper\_openings.json}。
\end{itemize}

测试结果位于：

\begin{Shaded}
\begin{Highlighting}[]
\NormalTok{dual\_v2\_1/11/book\_thickness\_000001\_plane\_corrected/}
\end{Highlighting}
\end{Shaded}

该测试证明侧面检测和安全拒绝分支已经运行，但没有每本书的卡尺真值，因此不能仅凭
该结果宣称达到毫米级绝对精度。

\paragraph{6.1.10
客户使用要求}\label{6110-ux5ba2ux6237ux4f7fux7528ux8981ux6c42}

\begin{enumerate}
\def\labelenumi{\arabic{enumi}.}
\tightlist
\item
  相机尽量正对书架，确保书脊完整、清晰、不过曝。
\item
  不要修改 RGB 或 depth 的尺寸、方向和裁剪区域。
\item
  避免强反光、透明塑封和大面积无纹理表面。
\item
  上线前至少抽查 5 本不同厚度的书并使用卡尺对比。
\item
  单帧结果只能作为夹爪预张开初值。
\item
  实际夹紧必须使用力控、力矩限制或电机电流阈值。
\item
  未经卡尺标定和人工复核，不得使用单帧结果直接执行无人值守抓取。
\end{enumerate}

\subsection{7.
查看边界与书本编号}\label{7-ux67e5ux770bux8fb9ux754cux4e0eux4e66ux672cux7f16ux53f7}

输出目录：

\begin{Shaded}
\begin{Highlighting}[]
\BuiltInTok{export} \VariableTok{RESULT}\OperatorTok{=}\VariableTok{$\{RESULT}\OperatorTok{:{-}}\VariableTok{$CAPTURE}\NormalTok{/book\_thickness\_analysis}\VariableTok{\}}
\end{Highlighting}
\end{Shaded}

全部书本边界：

\begin{Shaded}
\begin{Highlighting}[]
\FunctionTok{xdg{-}open} \StringTok{"}\VariableTok{$RESULT}\StringTok{/books\_thickness\_overlay.png"}
\end{Highlighting}
\end{Shaded}

分别查看三层书架：

\begin{Shaded}
\begin{Highlighting}[]
\FunctionTok{xdg{-}open} \StringTok{"}\VariableTok{$RESULT}\StringTok{/shelf\_1\_thickness.png"}
\FunctionTok{xdg{-}open} \StringTok{"}\VariableTok{$RESULT}\StringTok{/shelf\_2\_thickness.png"}
\FunctionTok{xdg{-}open} \StringTok{"}\VariableTok{$RESULT}\StringTok{/shelf\_3\_thickness.png"}
\end{Highlighting}
\end{Shaded}

书本编号规则：

\begin{itemize}
\tightlist
\item
  \texttt{S1-01}：第 1 层，从左向右第 1 本。
\item
  \texttt{S2-08}：第 2 层，从左向右第 8 本。
\item
  \texttt{S3-34}：第 3 层，从左向右第 34 本。
\end{itemize}

查看指定书本的独立核验图：

\begin{Shaded}
\begin{Highlighting}[]
\FunctionTok{xdg{-}open} \StringTok{"}\VariableTok{$RESULT}\StringTok{/book\_crops/S1{-}05.png"}
\end{Highlighting}
\end{Shaded}

图中：

\begin{itemize}
\tightlist
\item
  青色、橙色或绿色轮廓表示书脊边界。
\item
  \texttt{T} 表示估计厚度。
\item
  \texttt{+/-} 表示 95\% 工程不确定度。
\item
  \texttt{OPEN} 表示建议夹爪预张开宽度。
\end{itemize}

\subsection{8.
查询指定书本的夹爪开口}\label{8-ux67e5ux8be2ux6307ux5b9aux4e66ux672cux7684ux5939ux722aux5f00ux53e3}

以下命令查询 \texttt{S1-05}：

\begin{Shaded}
\begin{Highlighting}[]
\VariableTok{BOOK\_ID}\OperatorTok{=}\NormalTok{S1{-}05 }\StringTok{"}\VariableTok{$PYTHON}\StringTok{"} \AttributeTok{{-}} \OperatorTok{\textless{}\textless{}\textquotesingle{}PY\textquotesingle{}}
\StringTok{import json}
\StringTok{import os}
\StringTok{from pathlib import Path}

\StringTok{result = Path(os.environ["RESULT"])}
\StringTok{book\_id = os.environ["BOOK\_ID"]}
\StringTok{data = json.loads((result / "gripper\_openings.json").read\_text())}
\StringTok{item = data["commands"][book\_id]}

\StringTok{print("书本 ID:", book\_id)}
\StringTok{print("厚度:", round(item["measured\_thickness\_mm"], 2), "mm")}
\StringTok{print("95\% 不确定度:", round(item["uncertainty\_95\_mm"], 2), "mm")}
\StringTok{print("建议夹爪开口:", round(item["opening\_mm"], 2), "mm")}
\StringTok{print("建议夹爪开口:", round(item["opening\_m"], 6), "m")}
\StringTok{print("质量:", item["quality"])}
\OperatorTok{PY}
\end{Highlighting}
\end{Shaded}

注意：\texttt{RESULT} 必须先导出，而不是只定义为普通 shell 变量：

\begin{Shaded}
\begin{Highlighting}[]
\BuiltInTok{export} \VariableTok{RESULT}\OperatorTok{=}\VariableTok{$CAPTURE}\NormalTok{/book\_thickness\_analysis}
\end{Highlighting}
\end{Shaded}

查询其他书本时修改 \texttt{BOOK\_ID}：

\begin{Shaded}
\begin{Highlighting}[]
\VariableTok{BOOK\_ID}\OperatorTok{=}\NormalTok{S3{-}12 }\StringTok{"}\VariableTok{$PYTHON}\StringTok{"} \AttributeTok{{-}} \OperatorTok{\textless{}\textless{}\textquotesingle{}PY\textquotesingle{}}
\StringTok{import json}
\StringTok{import os}
\StringTok{from pathlib import Path}

\StringTok{data = json.loads(}
\StringTok{    (Path(os.environ["RESULT"]) / "gripper\_openings.json").read\_text()}
\StringTok{)}
\StringTok{item = data["commands"][os.environ["BOOK\_ID"]]}
\StringTok{print(item["opening\_mm"])}
\OperatorTok{PY}
\end{Highlighting}
\end{Shaded}

该命令只输出毫米开口值，方便传给机器人程序。

\subsection{9.
在机器人程序中读取}\label{9-ux5728ux673aux5668ux4ebaux7a0bux5e8fux4e2dux8bfbux53d6}

Python 示例：

\begin{Shaded}
\begin{Highlighting}[]
\ImportTok{import}\NormalTok{ json}
\ImportTok{from}\NormalTok{ pathlib }\ImportTok{import}\NormalTok{ Path}

\NormalTok{result }\OperatorTok{=}\NormalTok{ Path(}
    \StringTok{"/mnt/ssd/user/V/feature{-}splatting{-}inria/"}
    \StringTok{"2026{-}09{-}01{-}{-}23{-}22{-}25/book\_thickness\_analysis"}
\NormalTok{)}
\NormalTok{book\_id }\OperatorTok{=} \StringTok{"S1{-}05"}

\NormalTok{data }\OperatorTok{=}\NormalTok{ json.loads((result }\OperatorTok{/} \StringTok{"gripper\_openings.json"}\NormalTok{).read\_text())}
\NormalTok{measurement }\OperatorTok{=}\NormalTok{ data[}\StringTok{"commands"}\NormalTok{][book\_id]}

\NormalTok{opening\_mm }\OperatorTok{=}\NormalTok{ measurement[}\StringTok{"opening\_mm"}\NormalTok{]}
\NormalTok{opening\_m }\OperatorTok{=}\NormalTok{ measurement[}\StringTok{"opening\_m"}\NormalTok{]}
\NormalTok{quality }\OperatorTok{=}\NormalTok{ measurement[}\StringTok{"quality"}\NormalTok{]}

\ControlFlowTok{if}\NormalTok{ quality }\OperatorTok{!=} \StringTok{"high"}\NormalTok{:}
    \ControlFlowTok{raise} \PreprocessorTok{RuntimeError}\NormalTok{(}\SpecialStringTok{f"}\SpecialCharTok{\{}\NormalTok{book\_id}\SpecialCharTok{\}}\SpecialStringTok{ 需要人工复核"}\NormalTok{)}

\BuiltInTok{print}\NormalTok{(}\SpecialStringTok{f"夹爪目标：}\SpecialCharTok{\{}\NormalTok{opening\_mm}\SpecialCharTok{:.2f\}}\SpecialStringTok{ mm / }\SpecialCharTok{\{}\NormalTok{opening\_m}\SpecialCharTok{:.6f\}}\SpecialStringTok{ m"}\NormalTok{)}
\end{Highlighting}
\end{Shaded}

请根据夹爪控制接口选择 \texttt{opening\_mm} 或
\texttt{opening\_m}，不要混用单位。

\subsection{10.
修改安全间隙}\label{10-ux4feeux6539ux5b89ux5168ux95f4ux9699}

如果夹爪定位误差较大，可将安全间隙从 3 mm 增加到 5 mm：

\begin{Shaded}
\begin{Highlighting}[]
\VariableTok{CUDA\_VISIBLE\_DEVICES}\OperatorTok{=}\NormalTok{0 }\DataTypeTok{\textbackslash{}}
\StringTok{"}\VariableTok{$PYTHON}\StringTok{"}\NormalTok{ book\_thickness\_pipeline.py }\StringTok{"}\VariableTok{$CAPTURE}\StringTok{"} \DataTypeTok{\textbackslash{}}
  \AttributeTok{{-}{-}gripper{-}clearance{-}mm}\NormalTok{ 5}
\end{Highlighting}
\end{Shaded}

该参数只影响建议开口，不改变测得的书本厚度。

建议：

\begin{itemize}
\tightlist
\item
  高精度、小回差夹爪：2--3 mm。
\item
  一般机器人抓取：3--5 mm。
\item
  标定不足或运动误差较大：5 mm 以上，并配合低速接近。
\end{itemize}

\subsection{11.
增加多视角帧数}\label{11-ux589eux52a0ux591aux89c6ux89d2ux5e27ux6570}

需要更高的多视角稳定性时：

\begin{Shaded}
\begin{Highlighting}[]
\VariableTok{CUDA\_VISIBLE\_DEVICES}\OperatorTok{=}\NormalTok{0 }\DataTypeTok{\textbackslash{}}
\StringTok{"}\VariableTok{$PYTHON}\StringTok{"}\NormalTok{ book\_thickness\_pipeline.py }\StringTok{"}\VariableTok{$CAPTURE}\StringTok{"} \DataTypeTok{\textbackslash{}}
  \AttributeTok{{-}{-}frames}\NormalTok{ 0,8,16,24,32,40,48,56,64,71 }\DataTypeTok{\textbackslash{}}
  \AttributeTok{{-}{-}reference{-}frame}\NormalTok{ 24}
\end{Highlighting}
\end{Shaded}

参数说明：

\begin{itemize}
\tightlist
\item
  \texttt{-\/-frames}：参与厚度融合的帧编号。
\item
  \texttt{-\/-reference-frame}：用于书本编号、边界图和独立裁剪图的参考帧。
\end{itemize}

参考帧应满足：

\begin{itemize}
\tightlist
\item
  三层书架尽量完整可见。
\item
  图像清晰，没有运动模糊。
\item
  相机尽量正对书脊。
\item
  左右边缘的书本没有被裁掉。
\end{itemize}

增加帧数通常能改善随机误差，但不能修复错误分割或严重遮挡。

\subsection{12. 指定其他 GPU}\label{12-ux6307ux5b9aux5176ux4ed6-gpu}

使用第 2 块 GPU：

\begin{Shaded}
\begin{Highlighting}[]
\VariableTok{CUDA\_VISIBLE\_DEVICES}\OperatorTok{=}\NormalTok{1 }\DataTypeTok{\textbackslash{}}
\StringTok{"}\VariableTok{$PYTHON}\StringTok{"}\NormalTok{ book\_thickness\_pipeline.py }\StringTok{"}\VariableTok{$CAPTURE}\StringTok{"}
\end{Highlighting}
\end{Shaded}

查看 GPU 编号：

\begin{Shaded}
\begin{Highlighting}[]
\ExtensionTok{nvidia{-}smi}
\end{Highlighting}
\end{Shaded}

\subsection{13.
指定自定义输出目录}\label{13-ux6307ux5b9aux81eaux5b9aux4e49ux8f93ux51faux76eeux5f55}

\begin{Shaded}
\begin{Highlighting}[]
\VariableTok{CUDA\_VISIBLE\_DEVICES}\OperatorTok{=}\NormalTok{0 }\DataTypeTok{\textbackslash{}}
\StringTok{"}\VariableTok{$PYTHON}\StringTok{"}\NormalTok{ book\_thickness\_pipeline.py }\StringTok{"}\VariableTok{$CAPTURE}\StringTok{"} \DataTypeTok{\textbackslash{}}
  \AttributeTok{{-}{-}output} \StringTok{"}\VariableTok{$CAPTURE}\StringTok{/customer\_results"}
\end{Highlighting}
\end{Shaded}

之后设置：

\begin{Shaded}
\begin{Highlighting}[]
\BuiltInTok{export} \VariableTok{RESULT}\OperatorTok{=}\VariableTok{$CAPTURE}\NormalTok{/customer\_results}
\end{Highlighting}
\end{Shaded}

\subsection{14.
输出文件说明}\label{14-ux8f93ux51faux6587ux4ef6ux8bf4ux660e}

\begin{Shaded}
\begin{Highlighting}[]
\NormalTok{book\_thickness\_analysis/}
\NormalTok{├── book\_thickness.json}
\NormalTok{├── book\_thickness.csv}
\NormalTok{├── gripper\_openings.json}
\NormalTok{├── books\_thickness\_overlay.png}
\NormalTok{├── rejected\_books\_overlay.png}
\NormalTok{├── shelf\_1\_thickness.png}
\NormalTok{├── shelf\_2\_thickness.png}
\NormalTok{├── shelf\_3\_thickness.png}
\NormalTok{├── reference\_depth\_map.png}
\NormalTok{├── reference\_book\_masks.npz}
\NormalTok{└── book\_crops/}
\NormalTok{    ├── S1{-}01.png}
\NormalTok{    ├── S1{-}02.png}
\NormalTok{    └── ...}
\end{Highlighting}
\end{Shaded}

\begin{itemize}
\tightlist
\item
  \texttt{gripper\_openings.json}：机器人最方便读取的夹爪开口文件。
\item
  \texttt{book\_thickness.csv}：逐书汇总，可使用 Excel 打开。
\item
  \texttt{book\_thickness.json}：包含完整边界、逐帧测量和诊断信息。
\item
  \texttt{books\_thickness\_overlay.png}：全部书本的边界和厚度。
\item
  \texttt{rejected\_books\_overlay.png}：因侧面、观察角或平面质量不足而拒绝的候选。
\item
  \texttt{book\_crops/}：每本书的独立检查图。
\item
  \texttt{reference\_book\_masks.npz}：每本书的高分辨率二值掩码。
\end{itemize}

\subsection{15.
结果质量检查}\label{15-ux7ed3ux679cux8d28ux91cfux68c0ux67e5}

执行：

\begin{Shaded}
\begin{Highlighting}[]
\VariableTok{RESULT}\OperatorTok{=}\StringTok{"}\VariableTok{$RESULT}\StringTok{"} \StringTok{"}\VariableTok{$PYTHON}\StringTok{"} \AttributeTok{{-}} \OperatorTok{\textless{}\textless{}\textquotesingle{}PY\textquotesingle{}}
\StringTok{import json}
\StringTok{import os}
\StringTok{from collections import Counter}
\StringTok{from pathlib import Path}

\StringTok{result = Path(os.environ["RESULT"])}
\StringTok{data = json.loads((result / "book\_thickness.json").read\_text())}
\StringTok{quality = Counter(book["quality"] for book in data["books"])}

\StringTok{print("书本总数:", data["book\_count"])}
\StringTok{print("拒绝候选:", data["reference\_rejection\_count"])}
\StringTok{print("质量统计:", dict(quality))}
\StringTok{print("需要复核:")}
\StringTok{for book in data["books"]:}
\StringTok{    if book["quality"] != "high":}
\StringTok{        print(}
\StringTok{            book["book\_id"],}
\StringTok{            "厚度",}
\StringTok{            round(book["thickness\_mm"], 2),}
\StringTok{            "mm，不确定度",}
\StringTok{            round(book["uncertainty\_95\_mm"], 2),}
\StringTok{            "mm",}
\StringTok{        )}
\OperatorTok{PY}
\end{Highlighting}
\end{Shaded}

\texttt{quality} 含义：

\begin{itemize}
\tightlist
\item
  \texttt{high}：多视角一致性和匹配质量良好。
\item
  \texttt{medium}：可以作为抓取初值，但建议查看独立裁剪图。
\item
  \texttt{low}：不应自动执行抓取，需要重新扫描或人工确认。
\end{itemize}

常见拒绝原因：

\begin{itemize}
\tightlist
\item
  \texttt{inconsistent\_spine\_plane\_normal}：该书平面方向与画面中多数书脊不一致。
\item
  \texttt{low\_spine\_plane\_inlier\_ratio}：书脊平面有效深度点占比不足。
\item
  \texttt{low\_spine\_plane\_coverage}：有效平面没有覆盖足够的书脊长度。
\item
  \texttt{grazing\_view\_angle}：观察角度过于倾斜。
\item
  \texttt{side\_face\_dominates\_mask}：SAM 掩码中侧面占比过高。
\item
  \texttt{side\_face\_boundary\_ambiguous}：检测到侧面，但无法稳定确定交界位置。
\item
  \texttt{side\_face\_boundary\_unresolved}：深度不足以拟合侧面裁剪边界。
\end{itemize}

出现拒绝时，应调整相机位置或增加视角，不应手工把候选重新加入夹爪命令。

\subsection{16. 常见问题}\label{16-ux5e38ux89c1ux95eeux9898}

\subsubsection{\texorpdfstring{提示
\texttt{No\ module\ named\ \textquotesingle{}liblzfse\textquotesingle{}}}{提示 No module named \textquotesingle liblzfse\textquotesingle{}}}\label{ux63d0ux793a-no-module-named-liblzfse}

只有 Record3D \texttt{.depth} 输入需要该依赖。执行：

\begin{Shaded}
\begin{Highlighting}[]
\StringTok{"}\VariableTok{$PYTHON}\StringTok{"} \AttributeTok{{-}m}\NormalTok{ pip install pyliblzfse}
\end{Highlighting}
\end{Shaded}

通用 PNG/TIFF/NPY/NPZ 深度输入不需要 LZFSE 解码。

\subsubsection{提示 RGB 与 depth
宽高比不一致}\label{ux63d0ux793a-rgb-ux4e0e-depth-ux5bbdux9ad8ux6bd4ux4e0dux4e00ux81f4}

该错误表示 depth 没有按 RGB
图像坐标系完成配准，或者其中一个文件被单独裁剪、旋转 或缩放。应由相机
SDK 先把 depth 对齐到 RGB，不能仅通过强制缩放绕过检查。

\subsubsection{厚度整体大约相差 1000
倍}\label{ux539aux5ea6ux6574ux4f53ux5927ux7ea6ux76f8ux5dee-1000-ux500d}

这通常是深度单位错误：

\begin{itemize}
\tightlist
\item
  浮点深度默认按米读取。
\item
  16 位整数深度默认按毫米读取。
\item
  其他单位使用 \texttt{-\/-depth-scale} 指定``米/数值''。
\end{itemize}

例如深度数值 10000 表示 1 m 时，比例应为：

\begin{Shaded}
\begin{Highlighting}[]
\ExtensionTok{{-}{-}depth{-}scale}\NormalTok{ 0.0001}
\end{Highlighting}
\end{Shaded}

\subsubsection{提示找不到 SAM3
权重}\label{ux63d0ux793aux627eux4e0dux5230-sam3-ux6743ux91cd}

运行时指定权重：

\begin{Shaded}
\begin{Highlighting}[]
\VariableTok{CUDA\_VISIBLE\_DEVICES}\OperatorTok{=}\NormalTok{0 }\DataTypeTok{\textbackslash{}}
\StringTok{"}\VariableTok{$PYTHON}\StringTok{"}\NormalTok{ book\_thickness\_pipeline.py }\StringTok{"}\VariableTok{$CAPTURE}\StringTok{"} \DataTypeTok{\textbackslash{}}
  \AttributeTok{{-}{-}checkpoint}\NormalTok{ /actual/path/to/sam3.pt}
\end{Highlighting}
\end{Shaded}

\subsubsection{GPU 显存不足}\label{gpu-ux663eux5b58ux4e0dux8db3}

先查看 GPU：

\begin{Shaded}
\begin{Highlighting}[]
\ExtensionTok{nvidia{-}smi}
\end{Highlighting}
\end{Shaded}

选择空闲 GPU，例如：

\begin{Shaded}
\begin{Highlighting}[]
\VariableTok{CUDA\_VISIBLE\_DEVICES}\OperatorTok{=}\NormalTok{1 }\DataTypeTok{\textbackslash{}}
\StringTok{"}\VariableTok{$PYTHON}\StringTok{"}\NormalTok{ book\_thickness\_pipeline.py }\StringTok{"}\VariableTok{$CAPTURE}\StringTok{"}
\end{Highlighting}
\end{Shaded}

\subsubsection{检测书本数量不正确}\label{ux68c0ux6d4bux4e66ux672cux6570ux91cfux4e0dux6b63ux786e}

先尝试更换清晰、完整的参考帧：

\begin{Shaded}
\begin{Highlighting}[]
\VariableTok{CUDA\_VISIBLE\_DEVICES}\OperatorTok{=}\NormalTok{0 }\DataTypeTok{\textbackslash{}}
\StringTok{"}\VariableTok{$PYTHON}\StringTok{"}\NormalTok{ book\_thickness\_pipeline.py }\StringTok{"}\VariableTok{$CAPTURE}\StringTok{"} \DataTypeTok{\textbackslash{}}
  \AttributeTok{{-}{-}reference{-}frame}\NormalTok{ 12}
\end{Highlighting}
\end{Shaded}

仍有漏检时可略微降低阈值：

\begin{Shaded}
\begin{Highlighting}[]
\VariableTok{CUDA\_VISIBLE\_DEVICES}\OperatorTok{=}\NormalTok{0 }\DataTypeTok{\textbackslash{}}
\StringTok{"}\VariableTok{$PYTHON}\StringTok{"}\NormalTok{ book\_thickness\_pipeline.py }\StringTok{"}\VariableTok{$CAPTURE}\StringTok{"} \DataTypeTok{\textbackslash{}}
  \AttributeTok{{-}{-}min{-}score}\NormalTok{ 0.45}
\end{Highlighting}
\end{Shaded}

降低阈值可能增加重复检测，因此每次调整后必须检查
\texttt{books\_thickness\_overlay.png}。

\subsection{17.
上线前安全要求}\label{17-ux4e0aux7ebfux524dux5b89ux5168ux8981ux6c42}

\begin{enumerate}
\def\labelenumi{\arabic{enumi}.}
\tightlist
\item
  随机选择至少 5 本不同厚度的书，用卡尺测量实际厚度。
\item
  对比 \texttt{book\_thickness.csv}，确认是否存在稳定系统偏差。
\item
  检查夹爪控制接口使用的单位是 mm 还是 m。
\item
  首次抓取使用低速模式。
\item
  夹爪闭合必须启用力控、力矩限制或电流阈值。
\item
  \texttt{medium} 和 \texttt{low} 质量结果不得直接进行无人值守抓取。
\item
  书本紧密排列、手指无法插入时，应先执行抽书或建立间隙动作。
\end{enumerate}

本系统提供抓取预张开宽度，不替代机器人碰撞检测、力控和安全控制。

\end{document}
